\documentclass[11pt,a4paper]{article}
\usepackage[T1]{fontenc}
\usepackage[utf8]{inputenc}
\usepackage{amsmath,amssymb,xcolor}
\usepackage[margin=25mm]{geometry}
\usepackage[colorlinks=true,urlcolor=blue]{hyperref}
\definecolor{nexorablue}{HTML}{2459B5}
\setlength{\parindent}{0pt}
\setlength{\parskip}{7pt}
\setlength{\emergencystretch}{2em}
\renewcommand{\familydefault}{\sfdefault}
\begin{document}

{\small\bfseries\color{nexorablue} NEXORA CPGE}\par
{\small MP / NEX-RED-05 / Corrigé}\par
{\LARGE\bfseries Réduire et inverser à l’aide d’un polynôme annulateur}\par
\vspace{0.25em}\hrule\vspace{0.6em}

\section*{Spectre et projecteurs}
Le polynôme annulateur $(X-2)(X+1)$ est scindé à racines simples, donc $A$ est diagonalisable. Les projecteurs sont $(A+I)/3$ sur $E_2$ et $(2I-A)/3$ sur $E_{-1}$. Si une racine n’est pas une valeur propre, le projecteur correspondant est nul.

\section*{Inverse}
On a $A(A-I)=2I=(A-I)A$. Ainsi $A^{-1}=(A-I)/2$.


\end{document}
